\chapter{Integral Calculus and Multiple Integrals}
\label{chap:integral-calculus}

\section{Double Integrals and Fubini's Theorem}
Double integrals generalize single-variable Riemann integration to two-dimensional planar regions $R \subset \mathbb{R}^2$.
\begin{equation}
\iint_R f(x,y) \, dA = \lim_{\max \Delta A_k \to 0} \sum_{k=1}^N f(x_k, y_k) \Delta A_k
\end{equation}

\begin{theorem}[Fubini's Theorem]
If $f(x,y)$ is continuous on the rectangular region $R = [a, b] \times [c, d]$:
\begin{equation}
\iint_R f(x,y) \, dA = \int_a^b \left( \int_c^d f(x,y) \, dy \right) dx = \int_c^d \left( \int_a^b f(x,y) \, dx \right) dy
\end{equation}
For general regions $R = \{(x,y) \mid a \le x \le b, \, y_1(x) \le y \le y_2(x)\}$:
\begin{equation}
\iint_R f(x,y) \, dA = \int_a^b \left( \int_{y_1(x)}^{y_2(x)} f(x,y) \, dy \right) dx
\end{equation}
\end{theorem}

\section{Change of Order of Integration}
When an integrand is non-integrable with respect to the initial inner variable, changing the order of integration from $dy\,dx$ to $dx\,dy$ (or vice versa) resolves the evaluation.
\textbf{Methodology}:
\begin{enumerate}
    \item Identify the boundary curves from the original integration limits.
    \item Sketch the planar integration domain $R$.
    \item Construct horizontal strips instead of vertical strips (or vice versa) to read the new variable bounds.
\end{enumerate}

\begin{example}[Change of Order of Integration]
Evaluate $I = \int_0^1 \int_x^1 \frac{\sin y}{y} \, dy \, dx$.

\textbf{Solution:}
The inner integral $\int \frac{\sin y}{y} dy$ cannot be evaluated in elementary terms.
The region of integration is bounded by $y = x$, $y = 1$, $x = 0$, $x = 1$.
Reversing the order by slicing with horizontal strips from $x = 0$ to $x = y$, where $y$ ranges from $0$ to $1$:
\begin{align*}
I &= \int_0^1 \left( \int_0^y \frac{\sin y}{y} \, dx \right) dy = \int_0^1 \frac{\sin y}{y} [x]_0^y \, dy \\
  &= \int_0^1 \frac{\sin y}{y} (y) \, dy = \int_0^1 \sin y \, dy = [-\cos y]_0^1 = 1 - \cos 1
\end{align*}
\end{example}

\section{Change of Variables in Double Integrals}
Under coordinate transformation $x = x(u,v)$, $y = y(u,v)$:
\begin{equation}
\iint_R f(x,y) \, dx \, dy = \iint_{R'} f(x(u,v), y(u,v)) \left| \frac{\partial(x,y)}{\partial(u,v)} \right| \, du \, dv
\end{equation}
\textbf{Polar Coordinates ($x = r\cos\theta, y = r\sin\theta$)}:
The Jacobian is:
\begin{equation}
J = \begin{vmatrix} \frac{\partial x}{\partial r} & \frac{\partial x}{\partial \theta} \\ \frac{\partial y}{\partial r} & \frac{\partial y}{\partial \theta} \end{vmatrix} = \begin{vmatrix} \cos\theta & -r\sin\theta \\ \sin\theta & r\cos\theta \end{vmatrix} = r(\cos^2\theta + \sin^2\theta) = r
\end{equation}
Hence, the differential area element transforms as:
\begin{equation}
dx \, dy = r \, dr \, d\theta
\end{equation}

\section{Applications: Area and Volume}
\begin{itemize}
    \item \textbf{Planar Area of Region $R$}:
    \begin{equation}
    A = \iint_R dA = \iint_R dx \, dy = \iint_{R'} r \, dr \, d\theta
    \end{equation}
    \item \textbf{Volume Under Surface $z = f(x,y) \ge 0$ above $R$}:
    \begin{equation}
    V = \iint_R f(x,y) \, dx \, dy
    \end{equation}
\end{itemize}

\section{Triple Integrals and Spatial Coordinate Systems}
For a bounded three-dimensional solid $V \subset \mathbb{R}^3$:
\begin{equation}
\iiint_V f(x,y,z) \, dV = \int_{x_1}^{x_2} \int_{y_1(x)}^{y_2(x)} \int_{z_1(x,y)}^{z_2(x,y)} f(x,y,z) \, dz \, dy \, dx
\end{equation}

\begin{theorem}[Curvilinear Spatial Coordinates]
\begin{enumerate}
    \item \textbf{Cylindrical Coordinates $(r, \theta, z)$}:
    \begin{equation}
    x = r\cos\theta, \quad y = r\sin\theta, \quad z = z, \quad dV = r \, dr \, d\theta \, dz
    \end{equation}
    \item \textbf{Spherical Coordinates $(\rho, \phi, \theta)$}:
    \begin{equation}
    x = \rho \sin\phi \cos\theta, \quad y = \rho \sin\phi \sin\theta, \quad z = \rho \cos\phi, \quad dV = \rho^2 \sin\phi \, d\rho \, d\phi \, d\theta
    \end{equation}
\end{enumerate}
\end{theorem}

\begin{example}[Volume of a Sphere via Spherical Coordinates]
Calculate the volume of a sphere of radius $R$.

\textbf{Solution:}
In spherical coordinates: $\rho \in [0, R]$, $\phi \in [0, \pi]$, $\theta \in [0, 2\pi]$.
\begin{align*}
V &= \int_0^{2\pi} \int_0^{\pi} \int_0^R \rho^2 \sin\phi \, d\rho \, d\phi \, d\theta \\
  &= \left( \int_0^{2\pi} d\theta \right) \left( \int_0^{\pi} \sin\phi \, d\phi \right) \left( \int_0^R \rho^2 \, d\rho \right) \\
  &= (2\pi) \cdot [-\cos\phi]_0^{\pi} \cdot \left[ \frac{\rho^3}{3} \right]_0^R = 2\pi \cdot 2 \cdot \frac{R^3}{3} = \frac{4}{3}\pi R^3
\end{align*}
\end{example}
